\BourbakiExercises{习题}

\begin{enumerate}
\def\labelenumi{\arabic{enumi})}
\tightlist
\item
  \begin{enumerate}
  \def\labelenumii{\alph{enumii})}
  \tightlist
  \item
    给出一个诺特模 M 的例子，使得 M 的每个非零自同态都是单射，并且存在 M 的非零且非双射的自同态。
  \end{enumerate}
\end{enumerate}

\begin{enumerate}
\def\labelenumi{\alph{enumi})}
\setcounter{enumi}{1}
\tightlist
\item
  给出一个阿廷模 M 的例子，使得 M 的每个非零自同态都是满射，并且存在 M 的非零且非双射的自同态。
\end{enumerate}

\begin{enumerate}
\def\labelenumi{\arabic{enumi})}
\setcounter{enumi}{1}
\item
  设 \(\mathrm{A}\) 是一个环，其中双边理想的每个递增序列 \((\mathfrak{a}_n)_{n\in \mathbf{N}}\) 都是平稳的（例如，左或右诺特环）。证明环 \(\mathrm{A}\) 的每个满射自同态都是双射。
\item
  给出一个环的例子，它有两个元素 u、v，使得 uv = 1 且 \(vu \neq 1\)。证明这样的环不是诺特子环的一个有向族的并。
\item
  设 M 是一个 A-模，p 与 q 是 M 的两个投影算子。
\end{enumerate}

\begin{enumerate}
\def\labelenumi{\alph{enumi})}
\item
  我们有 \(p(\mathrm{M}) \subset q(\mathrm{M})\) 当且仅当 qp = p；核 \(\operatorname{Ker}(p)\) 包含于 \(\operatorname{Ker}(q)\) 当且仅当 qp = q。
\item
  和 \(p + q\) 是投影算子当且仅当 pq = -qp。若 \(p + q\) 是投影算子，且在 M 中关系 2x = 0 蕴含 x = 0，则我们有 pq = 0 与 qp = 0。
\item
  证明下列性质等价：
\end{enumerate}

\begin{enumerate}
\def\labelenumi{(\roman{enumi})}
\item
  我们有 \(pq = qp\)。
\item
  我们有 \(q(p(\mathbb{M})) \subset p(\mathbb{M})\) 与 \(q(\operatorname{Ker}(p)) \subset \operatorname{Ker}(p)\) 。
\item
  存在 M 的两两正交的投影算子的序列 \((p_1, p_2, p_3, p_4)\)，使得 \(p = p_1 + p_2\)、\(q = p_1 + p_3\)，且 \(1_{\mathrm{M}} = p_1 + p_2 + p_3 + p_4\) 。
\end{enumerate}

这样的序列（当它存在时）是唯一的：我们有 \(p_{1} = pq\) 。

\begin{enumerate}
\def\labelenumi{\arabic{enumi})}
\setcounter{enumi}{4}
\tightlist
\item
  设 \(\mathrm{A}\) 是一个环。
\end{enumerate}

\begin{enumerate}
\def\labelenumi{\alph{enumi})}
\item
  A 的元素 e 是幂等的当且仅当映射 \(x \mapsto xe\) 是左 A-模 \(A_{s}\) 的一个投影算子（相应地，映射 \(x \mapsto ex\) 是右 A-模 \(A_{d}\) 的一个投影算子）。
\item
  设 \(e_{1}, e_{2}\) 是环 A 中的幂等元。证明群 \(\mathrm{Hom}_{\mathrm{A}}(\mathrm{Ae}_{1}, \mathrm{Ae}_{2})\) 同构于 \(e_{1}A \cap Ae_{2} = e_{1}Ae_{2}\)，且环 \(\mathrm{End}_{\mathrm{A}}(\mathrm{Ae}_{1})\) 同构于环 \(e_{1}Ae_{1}\) 。
\item
  设 \(e_1, e_2\) 是 A 中的幂等元。证明下列性质等价：
\end{enumerate}

\begin{enumerate}
\def\labelenumi{(\roman{enumi})}
\item
  \(Ae_{1} = Ae_{2}.\)
\item
  \(e_{1}e_{2}=e_{1}\) 与 \(e_{2}e_{1}=e_{2}\) 。
\item
  \((1 - e_{1})\mathrm{A} = (1 - e_{2})\mathrm{A}.\)
\end{enumerate}

当这些性质成立时，存在可逆元 \(a \in \mathrm{A}\)，使 \(e_2 = ae_1a^{-1}\) 。

\begin{enumerate}
\def\labelenumi{\alph{enumi})}
\setcounter{enumi}{3}
\tightlist
\item
  设 \(e_{1}, e_{2}\) 是 A 中的两个幂等元。证明下列性质等价：
\end{enumerate}

\begin{enumerate}
\def\labelenumi{(\roman{enumi})}
\item
  左 A-模 \(Ae_{1}\) 与 \(Ae_{2}\) 同构。
\item
  右 A-模 \(e_1\mathrm{A}\) 与 \(e_2\mathrm{A}\) 同构。
\item
  存在元素 \(x, y\) 属于 \(A\)，使得 \(xy = e_1\) 且 \(yx = e_2\) 。
\end{enumerate}

当这些性质成立时，我们称 \(e_{1}\) 与 \(e_{2}\) 是等价幂等元。

\begin{enumerate}
\def\labelenumi{\alph{enumi})}
\setcounter{enumi}{4}
\tightlist
\item
  若 A 中的幂等元等价且属于 A 的中心，则它们相等。
\end{enumerate}

\(f\) ) 设 \((e_i)_{i\in \mathrm{I}}\) 与 \((f_{i})_{i\in \mathrm{I}}\) 是 A 中幂等元的有限族，使得 \(e_i e_j = 0\) 与 \(f_{i} f_{j} = 0\)（对 \(i\in \mathrm{I}, j\in \mathrm{I}, i\neq j\)），且 \(e_i\) 等价于 \(f_{i}\)（对每个 \(i\in \mathrm{I}\)）。则 \(e = \sum_{i\in \mathrm{I}}e_i\) 与 \(f = \sum_{i\in \mathrm{I}}f_{i}\) 是 A 中的等价幂等元。

¶ 6) 设 A 是环，u 与 v 是 A 的两个元素，使得 uv = 1 且 vu ≠ 1。对 i ≥ 1 与 j ≥ 1，令 \(e_{ij} = v^{i-1}u^{j-1} - v^i u^j\) 。

\begin{enumerate}
\def\labelenumi{\alph{enumi})}
\item
  证明我们有 \(e_{ij} \neq 0\) 与 \(e_{ij}e_{hk} = \delta_{jh}e_{ik}\)，其中 \(i,j,k,h\) 为整数 \(\geqslant 1\) 。
\item
  证明 A 的左理想 \(Ae_{ii}\)（对 \(i \geqslant 1\)）都非零、作为 A-模两两同构（利用习题 5 d)），且它们的和是直和。
\item
  证明我们有 \(u(v + e_{1i}) = 1\)（对每个 \(i \geqslant 1\)），且 \(e_{1i} \neq e_{1j}\)（对 \(i \geqslant 1, j \geqslant 1, i \neq j\)）。
\item
  设 B 是一个环，它具有 VIII，第 15 页习题 7 a) 的诸等价性质。由上题推出：B 的每个左可逆或右可逆的元素都是可逆的。
\end{enumerate}

\begin{enumerate}
\def\labelenumi{\arabic{enumi})}
\setcounter{enumi}{6}
\tightlist
\item
  用 A（相应地 a）表示形如 \(\mathrm{P}/(\mathrm{X}^{2}+1)^{n}\) 的有理分式组成的集合，其中 n 是整数 \(\geqslant0\)，\(P\in \mathbf{Q}[X]\) 是次数 \(\leqslant2n\)（相应地 \(\leqslant2n-1\)）的多项式。a) 证明 A 是体 \(\mathbf{Q}(X)\) 的一个子环，且 a 是环 A 的一个理想。b) A-模 \(a\oplus a\) 与 \(A\oplus A\) 同构。A-模 a 是投射的且有限生成的，但不同构于 A。
\end{enumerate}

\begin{enumerate}
\def\labelenumi{\alph{enumi})}
\setcounter{enumi}{2}
\tightlist
\item
  存在 \(\mathbf{K}\)，它是体 \(\mathbf{Q}\) 的一个二次扩张，使得 \(\mathrm{A}_{(\mathrm{K})}\) -模 \(\mathfrak{a}_{(\mathrm{K})}\) 与 \(\mathrm{A}_{(\mathrm{K})}\) 同构。
\end{enumerate}

\begin{enumerate}
\def\labelenumi{\arabic{enumi})}
\setcounter{enumi}{7}
\item
  设 A 是一个左阿廷环，并设 M 与 N 是 A-模 \(A_{s}^{n}\) 中两个互补的子模。若 M 是自由 A-模，则 N 是自由 A-模。
\item
  设 A 是一个完备拓扑环，m 是 A 的一个双边理想。假设环 A/m 是体，m 的每个元素都拓扑幂零（Gen.~Top., III, §6, 第 317 页，习题 9），且存在由 A 的加法群的子群组成的、A 中 0 的邻域的基本系。则环 A 是局部环，且 m 是其极大理想。特别地，对每个素数 p，p 进整数环 \(\mathbf{Z}_{p}\)（V, §12, No.~3, 第 96 页）是局部的，其极大理想是 \(p\mathbf{Z}_{p}\) 。
\end{enumerate}

¶ 10) 设 p 是素数，q 是 p 的幂，G 是 q 阶有限群，A 是交换环。设 \((e_{g})\) 是 A 上群 G 的代数 A{[}G{]} 的典范基。

\begin{enumerate}
\def\labelenumi{\alph{enumi})}
\item
  把 \(\mathfrak{a}\) 定义为元素 \(\sum a_{g}e_{g}\)（A{[}G{]} 中满足 \(\sum a_{g} = 0\) 者）组成的集合，它是 A{[}G{]} 的一个双边理想。若环 A 的特征为 \(p\)（V, §1, No.~2, 第 2 页），则我们有 \(x^{q} = 0\) 对每个 \(x\in \mathfrak{a}\) 成立（对 \(q\) 用归纳法；若 \(q\neq 1\)，取 G 的中心中一个元素 \(h\)，其阶为 \(p\)（I, §6, No.~5, 第 77 页，命题 11 的推论），并由归纳假设推出存在 \(y\in \mathrm{A}[\mathrm{G}]\) 使 \(x^{q / p} = (1 - e_h)y\)）。
\item
  设 A 是局部环，其极大理想为 m，且体 A/m 的特征为 p。证明环 A{[}G{]} 是局部环，且其极大理想是元素 \(\sum a_{g}e_{g}\in A[G]\)（满足 \(\sum a_{g}\in m\)）组成的集合。
\end{enumerate}

（设 S 是一个单 A{[}G{]}-模；证明 S 被 m 零化，然后，把 I, §6, No.~5, 第 76 页的命题 11 应用于由 S 的一个非零元素生成的 Z{[}G{]}-子模，证明 S 含有一个在 G 下不变的元素。）

\begin{enumerate}
\def\labelenumi{\arabic{enumi})}
\setcounter{enumi}{10}
\tightlist
\item
  设 p 是素数，G 是一个不含 p 阶元素的群，A 是一个不含非零幂零元素的、特征为 p 的交换环，A{[}G{]} 是 G 在 A 上的代数，\((e_{g})\) 是 A{[}G{]} 在 A 上的典范基。证明环 A{[}G{]} 没有非零的诣零理想（注意：若元素 \(\sum a_{g}e_{g} \in A[G]\) 幂零，则我们有 \(a_{1} = 0\)）。
\end{enumerate}

¶ 12) 设 K 是交换域，E 与 F 是 K 的扩张，A 是环 \(E \otimes_{K} F\) 。

\begin{enumerate}
\def\labelenumi{\alph{enumi})}
\tightlist
\item
  若 E 与 F 是 K 的超越扩张，则环 A 不是局部环。b) 设 E 是 K 的一个代数扩张，并用 \(E_{s}\) 与 \(F_{s}\) 分别表示 K 在 E 与 F 中的可分闭包。下列性质等价：
\end{enumerate}

\begin{enumerate}
\def\labelenumi{(\roman{enumi})}
\item
  环 A 是局部环。
\item
  环 \(\mathrm{E}\otimes_{\mathrm{K}}\mathrm{F}_s\) 是整环。
\item
  环 \(E \otimes_{K} F_{s}\) 是体。
\item
  环 \(E_{s} \otimes_{K} F_{s}\) 是体。
\end{enumerate}

当这些性质成立时，A 的极大理想 \(\mathfrak{m}\) 是 A 的幂零元素组成的集合，且 E 与 F 的每个复合扩张都同构于 \(A / \mathfrak{m}\) 。

\begin{enumerate}
\def\labelenumi{\arabic{enumi})}
\setcounter{enumi}{12}
\item
  设 A 是一个不是体的局部整环（例如 \(\mathbf{Z}_{(p)}\)，参见 VIII，第 26 页，例 5），并设 a 是 A 的一个非零不可逆元。把 A-模 \(\mathrm{A}^{(\mathbf{N})}\) 记作 M，其典范基记作 \((e_{n})_{n\in\mathbf{N}}\)。证明序列 \((f_{n})_{n\in\mathbf{N}}\)——其中 \(f_{2m}=e_{2m}\) 与 \(f_{2m+1}=e_{2m+1}+ae_{2m+2}\)（对每个 \(m\geqslant0\)）——是 M 在 A 上的一个基。证明 M 的子模 \(N=\sum_{m\geqslant0}\mathrm{A}(e_{2m}+ae_{2m+1})\) 是直因子，但不存在 N 的子集 J 使 \(\sum_{n\in J}Af_{j}\) 与 N 互补。由此推出：在 VIII，第 35 页的推论 5 中，不能略去 I 有限的假设。
\item
  设 A 是环，M 是 A-模，N 是 M 的一个原始直因子。设 \(M = \bigoplus_{i \in I} M_i\) 是 M 的一个直和分解，\((p_i)_{i \in I}\) 是相伴的投影算子族。证明存在 \(i \in I\)，使 \(p_i\) 诱导出从 N 到 \(M_i\) 的一个直因子的同构。
\item
  设 A 是环，并设 M、N、P 是 A-模。若模 \(M \oplus P\) 与 \(N \oplus P\) 同构且模 P 是原始的，则模 M 与 N 同构。
\item
  设 A 是环，M 是一个有限长度的 A-模，\((\mathrm{M}_{i})_{i\in\mathrm{I}}\) 是以 M 为直和的 M 的不可分解子模族。设 L 是一个有限长度的不可分解 A-模，\(I_{L}\) 是指标 \(i\in I\) 中使 \(M_{i}\) 同构于 L 的那些指标的集合。用例子说明 M 的子模 \(\sum_{i\in I_{L}}M_{i}\) 未必在 M 的自同构下稳定。
\item
  设 \(\mathrm{A}\) 是环，\(n\) 是整数，\(X \in \mathbf{GL}_n(\mathrm{A})\) 。
\end{enumerate}

\begin{enumerate}
\def\labelenumi{\alph{enumi})}
\item
  设环 A 是局部环。证明存在 \(\mathbf{GL}_{n}(\mathrm{A})\) 中的矩阵 P、Q，它们是形如 \(B_{ij}(\lambda)\) 的矩阵（II, §10, No.~13, 第 361 页）的乘积，以及 A 的一个可逆元 \(\delta\)，使得我们有 \(PXQ = \text{diag}(1, \ldots, 1, \delta)\)（照出处同前中命题 14 的证明进行）。
\item
  设环 A 是欧几里得环（VII, §1, 第 49 页，习题 7）。证明同样的结果。
\item
  若 \(\mathrm{A}\) 是局部环或欧几里得交换环，则群 \(\mathbf{S}\mathbf{L}_n(\mathbf{A})\) 由矩阵 \(B_{ij}(\lambda)\) 生成。
\item
  设 A 是局部环，其极大理想记作 m。证明 \(\mathbf{M}_{n}(\mathrm{A})\) 中每个在 \(\mathbf{M}_{n}(\mathrm{A}/\mathfrak{m})\) 中的像可逆的矩阵本身可逆（利用 II, §10, No.~13, 第 361 页的引理 1）。
\end{enumerate}

\begin{enumerate}
\def\labelenumi{\arabic{enumi})}
\setcounter{enumi}{17}
\tightlist
\item
  设 A 是局部环，E 是自由 A-模，x 是 E 的一个非零元素。设 \((e_{i})_{i\in I}\) 是 E 的一个使 x 的非零坐标个数尽可能少的基。设 \(x=\sum a_{i}e_{i}\) 是 x 在此基下的表示；用 J 表示指标 \(i\in I\) 中使 \(a_{i}\neq0\) 的那些指标的集合。
\end{enumerate}

\begin{enumerate}
\def\labelenumi{\alph{enumi})}
\tightlist
\item
  证明诸 \(a_{i}\) 中没有一个属于由其余元素生成的 A 的右理想。
\item
  设 F 与 G 是 E 的两个互补的子模，使得 \(x \in F\)；对 \(i \in I\)，写 \(e_{i} = f_{i} + g_{i}\)，其中 \(f_{i} \in F\)、\(g_{i} \in G\)。证明由 \(f_{j}\)（\(j \in J\)）与 \(e_{i}\)（\(i \in I - J\)）组成的族是 E 的一个基（若 \(g_{j} = \sum_{k \in I} b_{jk} e_{k}\)，注意 j、k 属于 J 时坐标 \(b_{jk}\) 属于 A 的极大理想，并应用习题 17 d)）。由此推出：存在 F 的一个自由子模，它包含 x 且是 F 的直因子。
\end{enumerate}

\(c\) ) 证明每个投射 A-模都是自由的（利用 Kaplansky 定理（II, §2, 第 386 页，习题 2）化为具有可数生成元族的模的情形，并应用 \(b\)））。

\begin{enumerate}
\def\labelenumi{\arabic{enumi})}
\setcounter{enumi}{18}
\item
  设 A 是交换环，M 是有限生成 A-模，u 是 M 的一个自同构，N 是 M 的一个在 u 下稳定的子模。证明我们有 \(u(N) = N\)（利用 VIII，第 28 页的推论 1, d)）。
\item
  设 A 是环，M 是一个有限长度的 A-模，N 是一个 A-模。设 m 与 n 是严格正整数，使得 \(M^{m}\) 与 \(N^{n}\) 同构。证明存在 A-模 P 与整数 p、q，使得 M 同构于 \(P^{p}\)，N 同构于 \(P^{q}\)，且 mp = nq。
\end{enumerate}

